%! Rational Functions and Zeros — Unit 1, Lesson 1.8 — Exit Ticket
\documentclass[10pt]{article}
\usepackage{precalculus-article}
\usepackage{precalculus-boxes}
\newcommand{\TallMath}[1]{$\displaystyle #1\rule[-1.4em]{0pt}{3.2em}$}

\begin{document}

\pageheader{Unit 1, Lesson 1.8}{Exit Ticket}
\namedateperiod

\vspace{0.14in}

\begin{enumerate}[itemsep=16pt]

\item Let $\displaystyle r(x) = \frac{x^2 - x - 6}{x + 4}$.

      \begin{enumerate}[label=(\alph*), itemsep=10pt]
        \item What value(s) of $x$ are \textbf{not} in the domain of $r$?
              Explain briefly.

              \vspace{0.08in}
              $x = $ \blank{3cm} because \blank{7cm}

        \item Find all real zeros of $r(x)$. Show your work.

              \vspace{0.06in}
              Factor the numerator: $x^2 - x - 6 = $ \blank{5cm}

              \vspace{0.08in}
              Numerator zeros: $x = $ \blank{4cm}

              \vspace{0.08in}
              Check each against the domain restriction: \blank{6cm}

              \vspace{0.08in}
              Zeros of $r$: $x = $ \blank{4cm}
      \end{enumerate}

\item A student says: ``$x = -4$ is a zero of
      $\displaystyle r(x) = \frac{x + 4}{x^2 - 1}$ because the numerator equals zero
      when $x = -4$.''

      Is the student correct? Circle one: \quad \textbf{Yes} \quad / \quad \textbf{No}

      \vspace{0.08in}
      Explain:

      \vspace{0.08in}
      \writelines{2}

\end{enumerate}

\end{document}
